% ---------------- Executive summary ----------------
\chapter*{Executive Summary}
\addcontentsline{toc}{chapter}{Executive Summary}
\thispagestyle{plain}

\textbf{The problem.} In Pakistan's board exams, millions of handwritten answer scripts are marked by examiners who are
paid per script, and only a small random sample is ever checked by a second person. Until recently, ``rechecking''
in Punjab only added up the marks again \citep{techjuice2025remark}. When a mark is wrong, a student usually has no real
way to get it corrected, and a few marks can decide a merit list, an admission, or a whole career path.

\textbf{The idea.} \emph{Parakh} is an audit layer that does not mark papers, but checks the marking.. For every scanned script it
reads the student's answers and the examiner's red marks, runs simple consistency rules, asks an AI ``second reader''
for an independent mark estimate, and ranks scripts by risk. Human reviewers then re-check the most suspicious scripts
first, instead of a random sample. The AI never gives the final mark.

\textbf{What I built and found.} A fully local prototype (one consumer GPU, open models, no cloud APIs), tested on a public
dataset of 50 real handwritten exam scripts \citep{dataset2026}:
\begin{itemize}
  \item Students' answers were read with about \textbf{99\%} accuracy; examiners' red numbers with only \textbf{27\%}. This is the main limitation.
  \item With the same 20\% re-checking effort, risk ranking caught \textbf{30\%} of planted marking errors vs \textbf{20\%} for random checking
        (average of 30 runs). A follow-up analysis showed \textbf{39\%} (about 2x random) when the unreliable red-ink check is removed.
  \item The dataset, described as cross-verified, already contained real marking errors. The audit ranked two of them
        \textbf{\#1 and \#2} of 50, and the AI second reader found a real \textbf{6-mark} under-marking.
  \item Estimated cost: about \textbf{\rs0.83 per script}, compared with about \rs35 paid for a human re-check (an estimate, see \dec{D021}).
\end{itemize}

\textbf{How to read this report.} Numbers marked with a decision number, e.g. \dec{D018}, can be traced to a file in the
project's \texttt{decisions/} folder, which records why each choice was made, which options were rejected, and what failed.
